\section{Literature}\label{sec:literature}

Three literatures meet here without previously having met each other:
the household incidence of trade policy, the microsimulation of shocks
through statutory tax--benefit systems, and the UK's episode-specific
empirical record since 2020.

\subsection{Trade to households}

The framework for pushing trade shocks onto household microdata is
best developed for developing countries: the World Bank's Household
Impacts of Tariffs programme simulates tariff changes through both
expenditure and income portfolios for 54 countries
\citep{artucporto2021, artucporto2019}, and the OECD has published the
concordance machinery for mapping CGE trade-model output onto
household-budget-survey categories \citep{luu2020}. The structural
anchor for the consumption side is \citet{fajgelbaumkhandelwal2016}:
trade's gains flow disproportionately to poorer households through
traded-goods budget shares. The 2025 tariff war produced a fast-moving
US applied literature---the Budget Lab's rolling distributional scoring
of the enacted schedules and the Tax Policy Center's implementation of
tariffs inside a working tax microsimulation model
\citep{budgetlab2025, tpc2025}---and a first academic treatment of
two-sided incidence on household expenditure microdata
\citep{fleckpradhan2026}. The CGE--microsimulation linkage tradition
supplies the architecture for putting a trade model upstairs and
microdata downstairs \citep{robilliardrobinson2005, peichl2008,
colombo2010, cali2019}; of that tradition, only \citet{peichl2008}
carries actual statutory tax rules in the micro layer.

What none of this work contains is a tax--benefit system responding to
the shock. The World Bank and OECD frameworks stop at market incomes
and expenditure; the US 2025 work is consumption-side; the
CGE-linkage papers (Peichl's flat-tax application aside) use
reduced-form income equations. As far as we can establish, no trade
shock has been run through a statutory tax--benefit calculator---
EUROMOD, UKMOD, TAXBEN, PolicyEngine or any counterpart---in any
country, for any episode.

\subsection{Shocks through statutory rules}

That machinery exists and has been applied to every kind of shock
except trade. \citet{dolls2012} measure the stabilisation coefficients
of the US and European systems against stylised income and unemployment
shocks; \citet{brewertasseva2021} impose the observed Covid
labour-market shock on FRS microdata through UKMOD and decompose the
cushioning into automatic and discretionary parts; the EUROMOD
nowcasting tradition supplies the labour-transition machinery
\citep{navicke2014}. The energy-price crisis produced the nearest
price-side analogues: distributional analyses of domestic energy costs
and the policy response by the IFS and Resolution Foundation, though
none frames the imported-price component as a trade shock or reports
stabilisation coefficients against it. The companion study
\citep{ahmadi2026} is, to our knowledge, the first tax--benefit
microsimulation of a trade shock anywhere---a single episode (the 2025
US tariffs), earnings channel only. This paper generalises it: four
episodes, both channels, and the rulebook vintage as an explicit
dimension.

\subsection{The UK episode record}

Each episode carries its own first-stage literature, reviewed with the
imported numbers in Section~\ref{sec:episodes}: the Brexit/TCA price
and trade effects \citep{bakker2023, breinlich2022, freeman2022}, the
2022 terms-of-trade and energy-price shock and its policy response, the
2025 US tariff schedule \citep{obr2025}, and the India agreement's
impact assessment \citep{dbtindia2026}. The distributional record of
these episodes is thin everywhere and absent at the household-rulebook
level: the Brexit price literature reports flat-to-regressive
consumer-cost incidence but never asks what the benefit system gave
back; the energy-crisis literature scores the discretionary response
but not the trade shock as such; the FTA impact assessments contain, in
DBT's own regulator's words, household analysis ``developed further for
future FTA IAs.'' The gap this paper fills is the join: the same
statutory second stage, run against all four first stages.
